\documentclass[10pt,a4paper]{amsart}
\usepackage[margin=19mm]{geometry}
\usepackage{amsmath,amssymb,mathtools}
\usepackage[scr=rsfs,cal=euler]{mathalfa}
\usepackage{xfrac}
\usepackage[unicode,hidelinks,bookmarks=false]{hyperref}
\newcommand{\ie}{i.e.\ }
\newcommand{\cf}{cf.\ }
\newcommand{\resp}{respectively\ }
\newcommand{\Subsection}{\subsection}
\providecommand{\nobreakdash}{\nobreak\hbox{-}}
\setlength{\emergencystretch}{2em}
\multlinegap0pt
\usepackage[shortlabels]{enumitem}
\usepackage{amssymb}
\usepackage{mathtools}
\usepackage{tikz}
\usepackage{dsfont}
\newtheorem{lemma}{Lemma}
\newtheorem{claim}{Claim}
\newtheorem{poc}{Poc}
\renewcommand\Pr{\operatorname{\mathbb{P}}}
\renewcommand\ge{\geqslant}
\renewcommand\le{\leqslant}
\title{Sharp threshold for Hamilton cycles in randomly perturbed sparse graphs}
\author{Guorui Ma, Zhifei Yan}
\date{Source: arXiv:2605.29553v2 (CC BY 4.0)}
\begin{document}
\maketitle

\noindent\emph{Source and licence.} Sharp threshold for Hamilton cycles in randomly perturbed sparse graphs. Version arXiv:2605.29553v2 (CC BY 4.0), distributed by arXiv under the Creative Commons Attribution 4.0 International licence, http://creativecommons.org/licenses/by/4.0/, as recorded in the arXiv licence metadata for this exact version and shown on the abstract page https://arxiv.org/abs/2605.29553v2. Full licence notice: \texttt{licenses/arxiv-2605.29553-CC-BY-4.0.txt}.\par\smallskip

\noindent\textit{Editorial scope.} Counted unit: the article's lemma lem:expansion (source0.tex lines 227--234). The article's complete original proof of it is delivered with it (source0.tex lines 236--275, one \texttt{proof} environment of 40 lines). No mathematics is written, repaired or re-derived: the statement and the proof are the article's own lines, in the article's own notation, and no retained line is substituted or re-flowed.  No further source-local context is needed: every cross-reference of every retained line resolves inside the two delivered spans.  Nothing was omitted from any retained span, so the package carries no omission record.  No retained line contains a citation, so the wrapper carries no bibliography and no bibliography entry.  No retained line contains a figure environment, an image-inclusion command or drawing code, so no image is delivered and the package declares no asset dependency.  Editorial formatting only, in addition: an AMS \texttt{amsart} preamble that restates the article's own declarations and macros used by the retained lines, namely the environments \texttt{lemma}, \texttt{claim}, \texttt{poc} on separate counters, together with the standard packages those lines need.  The retained environments print in this document as Lemma 1, Claim 1, Poc 1 in order of appearance, whereas the article prints the same items with the numbering of its own full document.  The complete native source of the article as distributed in the arXiv e-print for this exact version is bundled unchanged as \texttt{sources/201613/source0.tex}.
\begin{lemma}[Uniform random expansion]\label{lem:expansion}
Fix $\eta>0$. Let $\alpha=o(1)$, $L=\log(1/\alpha)$, $d=\lceil\alpha n\rceil$, and $\lambda=(1+\eta)L$. Assume $d\gg \log n$. There is a constant $K=K(\eta)>0$ such that $R\sim G(n,\lambda/n)$ has the following three properties a.a.s.
\begin{enumerate}[label=\textnormal{(E\arabic*)}]
    \item \label{E1} For every set $X$ with $\frac{Kd}{\lambda}\le|X|\le\frac{Kn}{2\lambda}$, we have $|N_{R}(X)\setminus X|>\frac{\lambda|X|}{K}$.
    \item \label{E2} For every set $X$ with $\frac{Kn}{2\lambda}<|X|\le\frac{n}{4}$, we have $|N_{R}(X)\setminus X|>\frac{n}{2}$.
    \item \label{E3} Every two disjoint sets $A, B\subseteq V$ with $|A|,|B|\ge n/4$ have at least one $R$-edge between them.
\end{enumerate}
\end{lemma}

\begin{proof}
We choose $K$ sufficiently large in terms of $\eta$. Let $X$ be fixed with $|X|=x$, and put $Z_{X}=|N_{R}(X)\setminus X|$. Then $Z_{X}\sim \text{Bin}(n-x,q_{x})$ where $q_{x}=1-(1-\lambda/n)^{x}$.

\begin{claim}\label{claim:e1}
Property \ref{E1} holds a.a.s.
\end{claim}
\begin{poc}
Write $a=\frac{\lambda x}{n}$. Then $K\alpha \le a\le K/2$. Choose a small constant $a_{0}>0$. If $a\le a_{0}$, then $\mathbb{E}Z_{X}=(1+O(a_{0})+o(1))an,$ and the standard binomial lower-tail estimate gives
$$\Pr\left(Z_{X}\le\frac{an}{K}\right)\le \exp\left[-an\left(1-O(a_{0})-\frac{1+\log K}{K}+o(1)\right)\right].$$
On the other hand,
$$\log\binom{n}{x}\le x \log\frac{en}{x}=\frac{an}{\lambda}\log\frac{e\lambda}{a}.$$
Since $a\ge K\alpha$ and $\lambda=(1+\eta)L,$ we have
$$\frac{1}{\lambda}\log\frac{e\lambda}{a}\le\frac{1+o(1)}{1+\eta}.$$
Choosing first $a_{0}$ sufficiently small and then $K$ sufficiently large, the union-bound exponent is at most $-can$ for some $c=c(\eta)>0$. Since $an=\lambda x\ge Kd$ and $d\gg \log n$, the total contribution from the subrange $a\le a_{0}$ is $o(1)$.

If $a\ge a_{0}$, then $\mathbb{E}Z_{X}=\Omega(n)$, while the threshold $an/K$ is a fixed smaller fraction of $\mathbb{E}Z_{X}$ once $K$ is chosen large. Thus $\Pr(Z_{X}\le an/K)\le \exp(-\Omega(n))$. The number of possible sets $X$ in the present range is $\exp(o(n))$, because $x=O(n/\lambda)$ and $\lambda\rightarrow\infty$. Hence this subrange also contributes $o(1)$. This proves \ref{E1}.
\end{poc}

\begin{claim}\label{claim:e2}
Property \ref{E2} holds a.a.s.
\end{claim}
\begin{poc}
If $Z_{X}\le n/2$ and $x\le n/4,$ then at least $n/4$ vertices outside $X$ have no $R$-neighbour in $X$. Hence, for fixed $X$,
$$\Pr(Z_{X}\le n/2)\le\binom{n}{n/4}\exp\left(-\frac{\lambda x}{4}\right).$$
Taking a union bound over all $X$ with $Kn/(2\lambda)<x\le n/4$ gives at most
$$n \exp\left(2H(1/4)n-\frac{Kn}{8}\right),$$
where $H(p)=-p\log p-(1-p)\log(1-p)$ is the binary entropy function. Taking $K>16H(1/4)$ proves \ref{E2}.
\end{poc}

\begin{claim}\label{claim:e3}
Property \ref{E3} holds a.a.s.
\end{claim}
\begin{poc}
For fixed disjoint $A, B$ with $|A|,|B|\ge n/4$, the probability of having no edges is
$$\Pr(e_{R}(A,B)=0)\le \exp(-\lambda n/16).$$
There are at most $4^{n}$ ordered choices of such pairs, and $\lambda\rightarrow\infty$. Thus the union bound proves \ref{E3} as required.
\end{poc}

The three claims complete the proof of the lemma.
\end{proof}

\end{document}
